\documentclass[a4paper,12pt]{article}
\usepackage{amsmath}
\usepackage[margin=1in]{geometry}
\begin{document}

\title {Dimensional Analysis and Rates of Change}
\maketitle
		

\section{Dimensional Analysis}

Energy has SI units of Joules, represented by a $J$. The Joule is derived from
base units as follows:

\centerline{$J$ $=$ $kg$ $m^2$ $s^{-2}$}
Let $M$ be mass, $A$ be acceleration, and $D$ be distance. Find a formula for Energy by writing the equation

\centerline{$E$ = $kg$ $m^2$ $s^{-2}$ $=$ $M^\alpha A^\beta D^\gamma$} then replacing A with the usual units for acceleration, D with the usual units for distance, and M with the usual units for mass. 
Now you have an equation where both sides have the same bases.
\begin{enumerate} 
	\item{Find the exponents $\alpha, \beta$, and $\gamma$.}
	\item{Knowing that $F=MA$ is the equation for Force, write the units of Force in terms of $kg, m$, and $s$}
	\item{If a ball is dropped with a mass 1kg and it travels a distance 1m to the floor, write an expression for the amount of energy it has when it hits the floor, including the units. Write an expression for the amount of force it has.}
	\item{Compare the units of Force and Energy. Dividing $\frac{E}{F}$, what are the units of the result? What is the relation between Force and Energy?}



\end{enumerate}

\section{Bonus question}

Here you will use dimensional analysis to derive a previously unseen force equation. For a radiating object, like a glowing mass of gas out in space, let the Intensity of the light be given by $I$ with units of $kg s^{-3}$. Let A be the area of the beam of light being measured, and c is the speed of light.
\begin{enumerate}
	\item{Write an equation for the radiation pressure force of the beam of light. Use dimensional analysis to do so, knowing that it must have the same units as force.}
\end{enumerate}
There is only one way to put the variables together to get the correct units. In this manner, even if you don't understand the Physics behind the question, you can still derive the correct formulae (up to a constant) and therefore nearly the write numerical answer. 
\end{document}
